\cleardoublepage % memastikan bab baru dimulai di halaman ganjil (kanan)
\chapter{PENDAHULUAN}
\label{chap:pendahuluan}
% --- Latar Belakang ---
\section{Latar Belakang}
\label{chap:pd_latbel}
Model bahasa besar atau \textit{Large Language Model} (LLM) yang awalnya diaplikasikan sebagai alat untuk memahami mekanisme otak manusia, kini mulai diaplikasikan ke arah sebaliknya, yaitu pemanfaatan sinyal otak manusia untuk memperbaiki model bahasa besar. Sinyal tersebut merupakan salah satu bagian penting dari objek penelitian \textit{neuroscience}, yaitu disiplin ilmu dengan fokus pada kajian sistem saraf otak, termasuk yang mengkaji mekanisme otak manusia dalam memproses bahasa. Pemahaman akan mekanisme tersebut berpotensi besar untuk diterapkan pada pengembangan model bahasa besar dengan representasi internal yang lebih selaras dengan cara kerja otak dalam memproses bahasa. Penerapan ini diharapkan dapat meningkatkan kapabilitas model dalam memahami konteks bahasa secara semantik.

Pengaplikasian LLM ke arah tersebut memanfaatkan data sinyal otak manusia seperti \textit{functional magnetic resonance imaging} atau fMRI. Secara khusus, \textcite{lebel2023natural} merilis data fMRI berskala besar berisi rekaman lengkap aktivitas otak saat sejumlah 8 partisipan mendengarkan siniar (\textit{podcast}). Ketersediaan data ini mendukung penelitian mendalam sebagaimana telah dijabarkan sebelumnya, terlebih data turunannya (untuk keperluan komputasi) sudah menerapkan prosedur \textit{data preprocessing}. Selain itu, data \textit{electrocorticography} (ECoG) dan \textit{magnetoencephalography} (MEG) juga digunakan pada penelitian serupa yang berfokus pada pemanfaatan sinyal otak manusia untuk meningkatkan pemahaman semantik model bahasa besar. 

Ketersediaan data sinyal otak mendorong eksplorasi \textit{brain-tuning} yang diaplikasikan pada tiga segmentasi. Pertama, pembuktian akan adanya keselarasan representasi internal antara jaringan otak manusia dengan LLM \autocite{goldstein2022shared, schrimpf2021neural, mischler2024contextual, yu2024predicting}. Kedua, perbaikan model berbasis bukti peningkatan pemahaman semantik atau kecerdasan LLM \autocite{schwartz2019inducing, moussa2026generalizability, moussa2025improving, negi2025brain, xiao2026beyond}. Ketiga, implementasi nyata sinyal otak untuk merekonstruksi teks \autocite{Tang2023858}.

% Beda proses implementasi -> Research gap -> Usulan gua

Fokus utama riset ini terletak pada pengaplikasian \textit{brain-tuning} pada segmen kedua, yaitu perbaikan LLM yang terukur berdasarkan peningkatan pemahaman semantik. Serangkaian tinjauan literatur pada segmen tersebut menggugah temuan akan adanya celah riset yang potensial. Dari segi pemanfaatan model, jenis model yang digunakan bervariasi pada tiga kelompok: model ucapan atau \textit{speech-LLM} \autocite{moussa2026generalizability, moussa2025improving} model teks dengan parameter kecil (<1 miliar) \autocite{schwartz2019inducing, negi2025brain}, serta model teks dengan parameter besar (1.5--72 miliar) untuk penalaran deduktif \autocite{xiao2026beyond}. Dari segi evaluasi, umumnya evaluasi berhenti pada skor akhir seperti kinerja tugas hilir atau korelasi \textit{encoding}. Dari segi galat riset, efek samping paparan fitur rendah seperti panjang dan frekuensi kata berpotensi menimbulkan bias dalam meningkatkan pemahaman semantik. Dari segi fungsionalitas model, \textcite{xiao2026beyond} mencatat bahwa kenaikan skor kecerdasan model tidak selalu sejalan dengan manfaat fungsionalitasnya. 

Berdasarkan tinjauan tersebut, pengaplikasian \textit{brain-tuning} pada riset ini dapat dikembangkan melalui beberapa pendekatan. (1) Memanfaatkan LLM teks berarsitektur \textit{decoder} dengan parameter realistis (1--8 miliar) yang diseimbangkan melalui mekanisme PEFT (\textit{Parameter-Efficient Fine-Tuning}) seperti LoRA (\textit{Low-Rank Adaptation}) menggunakan data fMRI naturalistik milik \textcite{lebel2023natural}. (2) Melakukan evaluasi melalui tugas pemahaman semantik hilir dengan variasi model pembanding berupa model yang di-\textit{fine-tuning} bahasa saja dan model yang dilakukan pengacakan (\textit{randomized}) sinyal otak. (3) Mengalisis secara kuantitatif representasi internal model (per lapisan) melalui tiga hal: perubahan per lapisan model dan per area otak (skor hierarki dan CKA), kontribusi fitur teks tingkat rendah (dekomposisi variansi), dan keseimbangan antara keselarasan otak, skor semantik, serta fungsionalitas model dengan memvariasikan bobot \textit{loss} sinyal otak. (4) Analisis pada tiap tingkat \textit{attention head} (misalnya penonaktifan LoRA per lapisan atau \textit{head}) yang dilakukan secara opsional apabila waktu riset memungkinkan. Pengembangan ini memberi manfaat ilmiah berupa pemetaan representasi internal LLM, serta validasi peran sinyal otak dalam memberi informasi yang tidak dapat diperoleh hanya dari teks saja. % ini nanti dikaji lagi brads..

Dari penjabaran tersebut, celah yang dituju riset ini adalah analisis internal model. Sejauh literatur yang ditinjau, belum ditemukan riset yang menganalisis secara kuantitatif bagaimana representasi per lapisan LLM teks berubah setelah \textit{brain-tuning} dengan fMRI naturalistik. Penelitian ini mengisi celah tersebut pada LLM teks berarsitektur \textit{decoder}. Selain mengukur peningkatan pemahaman semantik, riset ini menjelaskan jenis dan alasan dari setiap perubahan yang terjadi. Rumusan masalah yang diturunkan dari celah riset tersebut dijabarkan secara merinci pada Subbab \ref{chap:pd_rm}.

% --- Rumusan Masalah ---
\section{Rumusan Masalah}
\label{chap:pd_rm}
Dalam pelaksanaannya untuk memperoleh hasil baik dalam bentuk pemerolehan model dengan peningkatan pemahaman semantik ataupun manfaat ilmiah dari informasi luaran riset, penting untuk memastikan beberapa hal fundamental. Hal-hal tersebut mencakup mekanisme perancangan riset \textit{brain-tuning} yang memanfaatkan beberapa sumber daya komputasional (model, data, dan pengaplikasian PEFT), validasi peningkatan pemahaman semantik model pada tugas hilir serta kontrol keseimbangan fungsionalitas model, dan pemahaman akan perubahan representasi internal yang berubah dengan mempertimbangkan efek samping fitur teks tingkat rendah. Cakupan tersebut dirumuskan secara komprehensif pada rumusan masalah berikut:

\begin{enumerate}[label=RM\arabic*.]
    \item Bagaimana merancang dan mengimplementasikan \textit{brain-tuning} berbasis pengaplikasian PEFT (LoRA) pada LLM teks berarsitektur \textit{decoder} (1--8 miliar parameter) dengan objektif \textit{brain-alignment} dari data fMRI naturalistik?
    \item Apakah \textit{brain-tuning} meningkatkan pemahaman semantik hilir dibandingkan \textit{fine-tuning} bahasa saja dan pengacakan sinyal otak, dan bagaimana keseimbangannya dengan keselarasan otak serta fungsionalitas model ketika bobot \textit{loss} dari sinyal otak manusia divariasikan nilainya?
    \item Bagaimana representasi internal LLM berubah per lapisan model dan per area otak akibat \textit{brain-tuning}, dan seberapa besar kenaikan keselarasan yang berasal dari fitur teks tingkat rendah dibandingkan informasi semantik?
\end{enumerate}

Setiap rumusan masalah akan dipetakan pada beberapa hipotesis di Subbab \ref{chap:pd_hipo}. Secara khusus, RM1 tidak akan dipetakan pada hipotesis karena bersifat perancangan.

% --- Tujuan ---
\section{Tujuan}
Berdasarkan ketiga rumusan masalah yang telah dijabarkan, riset ini mengharapkan ketercapaian luaran yang dapat diukur. Ketercapaian tersebut disusun menjadi beberapa tujuan tugas akhir yang dipetakan dari tiap poin rumusan masalah berikut:
\begin{enumerate}
    \item Mengimplementasikan \textit{brain-tuning} berbasis LoRA pada LLM teks \textit{decoder} (1--8 miliar parameter) dengan objektif \textit{brain-alignment} dari fMRI naturalistik, termasuk menentukan lapisan target, objektif, dan bobot \textit{loss} dasar.
    \item Mengevaluasi pengaruh \textit{brain-tuning} terhadap pemahaman semantik hilir dibandingkan \textit{fine-tuning} bahasa saja dan pengacakan sinyal otak, serta membuat pemetaan keseimbangan (\textit{balance-mapping}) antara keselarasan otak, skor semantik, dan fungsionalitas model terhadap bobot \textit{loss}.
    \item Menganalisis secara kuantitatif akan perubahan representasi internal per lapisan model dan per area otak, serta mengukur kontribusi fitur teks tingkat rendah terhadap peningkatan keselarasan.
\end{enumerate}

% --- Batasan Masalah ---
\section{Ruang Lingkup, Asumsi, dan Batasan Masalah}
Keberjalanan tugas akhir ini akan tetap mempertahankan rasionalitas dan terlaksana secara realistis. Dalam upaya mempertahankan hal tersebut, berikut adalah rincian ruang lingkup, asumsi, serta batasan masalah dari seluruh pelaksanaan tugas akhir:

\subsection*{Ruang Lingkup}
\begin{enumerate}
    \item Perancangan \textit{brain-tuning} berfokus pada LLM teks \textit{decoder} berukuran 1--8 miliar parameter yaitu: Llama-3-8B LLama3.1-8B, Qwen3-4B, dan/atau Mistral-7B (bersifat opsional) menggunakan LoRA pada rentang 1--4 miliar parameter dan QLoRA (bersifat opsional) pada rentang 4--8 miliar parameter.
    \item Data fMRI naturalistik yang digunakan untuk \textit{brain-tuning} bersumber dari dataset publik \textcite{lebel2023natural}, dengan dataset Narratives milik \textcite{nastase2021narratives} (bersifat opsional) yang memvalidasi generalisasi lintas dataset.
    \item Evaluasi tugas hilir (\textit{downstream task}) pada pemahaman semantik dengan \textit{benchmark}: STS-B (kemiripan makna kalimat), WiC/SuperGLUE (makna kata dalam konteks), RTE atau MNLI (inferensi bahasa alami), dan HellaSwag (kelanjutan narasi). Pada fungsionalitas model dengan \textit{benchmark}: MMLU (pengetahuan dan juga penalaran) dan Perplexity (kemampuan bahasa dasar). 
\end{enumerate}

\subsection*{Asumsi}
\begin{enumerate}
    \item Respon BOLD dapat dipetakan secara linear dari representasi internal model.
    \item Tahap prapemrosesan pada dataset turunan yang tersedia valid dan konsisten.
    \item Penundaan hemodinamik dimodelkan dengan penundaan tetap beberapa TR.
    \item Penjajaran waktu untuk setiap kata pada transkrip dengan sinyal fMRI akurat.
    \item \textit{Benchmark} yang dipilih mengukur pemahaman semantik dan \textit{probing} linear pada representasi beku mencerminkan informasi yang dapat dikodekan LLM.
    \item Pengacakan TR (\textit{Repetition Time}) menghilangkan informasi spesifik otak tanpa mengubah statistik dasar sinyal otak manusia pada dataset original fMRI.
    \item Transkrip berpenanda waktu cukup merepresentasikan stimulus untuk menjawab rumusan masalah, tanpa adanya tambahan dari informasi akustik dataset.
    \item Luaran model terpilih cukup mewakili LLM pada rentang ukuran yang sama.
\end{enumerate}

\subsection*{Batasan Masalah}
\begin{enumerate}
    \item Jenis LLM terbatas pada teks \textit{decoder}. Riset tidak mencakup jenis LLM \textit{speech}, \textit{encoder}, atau multimodal, sehingga hasil tidak berlaku untuk semua LLM.
    \item Data sinyal otak yang digunakan hanyalah fMRI, dan tidak menggunakan ECoG, MEG, EEG, atau data serupa lainnya. Data Narratives digunakan secara opsional (hanya jika waktu memungkinkan) pada evaluasi lintas dataset.
    \item Bahasa yang digunakan baik pada data fMRI dan setiap \textit{benchmark} adalah bahasa Inggris (partisipan mendengar siniar \textit{Moth Radio Hours} pada perekaman fMRI), sehingga luaran yang dihasilkan tidak digeneralisasi pada bahasa lain.
    \item Jumlah partisipan pada dataset fMRI terbatas hanya pada 8 orang partisipan, sehingga luaran yang dihasilkan tidak diklaim berlaku untuk populasi umum.
    \item Masukan model hanya berupa transkrip teks, sehingga informasi akustik dan prosodi tidak disertakan dalam pemodelan. Selain itu, area auditori primer pada otak manusia juga tidak diharapkan sepenuhnya dijelaskan oleh model.
    \item Analisis tingkat sirkuit (\textit{causal tracing}) tidak dilakukan karena ada di luar cakupan riset. Analisis per \textit{attention head} bersifat opsional jika memungkinkan.
    \item Perbandingan dengan kontrol hanya menunjukkan selisih, tidak membuktikan adanya hubungan kausal sinyal otak manusia terhadap kemampuan bahasa.
\end{enumerate}

% --- Hipotesis ---
\section{Hipotugas}
\label{chap:pd_hipo}
Hipotesis pada tugas akhir ini disusun berdasarkan premis-premis yang diperoleh dari latar belakang beserta rumusan masalah. Berikut tersaji premis tugas akhir ini:

\subsection*{Premis}
\begin{enumerate}[label=P\arabic*.]
    \item Representasi internal LLM selaras sebagian dengan aktivitas otak manusia saat memproses bahasa natural \autocite{goldstein2022shared, schrimpf2021neural}.
    \item Hierarki lapisan model sejajar dengan hierarki korteks: lapisan lebih dalam selaras dengan area otak yang jauh dari korteks auditori \autocite{mischler2024contextual}.
    \item Pada penalaran deduktif, proses \textit{brain-tuning} dengan sinyal otak (fMRI) asli memberi hasil lebih baik daripada sinyal otak yang diacak \autocite{xiao2026beyond}.
    \item Kenaikan keselarasan global (dalam skor kecerdasan) tidak selalu sejalan dengan fungsionalitas LLM. Selain itu, fungsionalitas LLM stabil pada bobot \textit{loss} sinyal otak yang relatif kecil \autocite{xiao2026beyond}. Secara logis, kemampuan tersebut berpotensi mengalami penurunan pada bobot \textit{loss} yang besar. Hal ini dapat disebabkan oleh adanya dominasi objektif otak pada objektif bahasa.
    \item Pada LLM ucapan (\textit{speech-LLM}), proses \textit{brain-tuning} terbukti dapat mengurangi pengaruh fitur tingkat rendah terhadap keselarasan \autocite{moussa2025improving}.
    \item Dampak \textit{brain-tuning} terhadap pemahaman semantik LLM teks dan perubahan representasi per lapisan perlu dianalisis secara kuantitatif \autocite{negi2025brain}.
\end{enumerate}

% \begin{enumerate}
%     \item Representasi internal model bahasa dapat digunakan untuk memprediksi aktivitas otak manusia saat memproses bahasa \autocite{schrimpf2021neural}.
%     \item \textit{Fine-tuning} model bahasa menggunakan sinyal otak (fMRI dan MEG) dapat dilakukan tanpa menimbulkan kerusakan pada kapabilitas model \autocite{schwartz2019inducing}.
%     \item \textit{Brain-tuning} pada model ucapan terbukti meningkatkan pemahaman semantik model \autocite{moussa2025improving}.
%     \item Dataset fMRI naturalistik berskala besar dengan rekaman puluhan jam per partisipan telah tersedia secara publik \autocite{lebel2023natural}.
%     \item Penelitian sebelumnya belum menerapkan \textit{brain-tuning} pada LLM modern berskala besar dengan kontrol ketat yang memisahkan kontribusi sinyal otak dari efek \textit{fine-tuning} lainnya.
% \end{enumerate}

\subsection*{Hipotesis}
Berdasarkan premis-premis tersebut, berikut tersaji hipotesis tugas akhir berupa pemetaan tiap premis dengan dugaan-dugaan awal proses riset yang terukur untuk pembuktian secara teoritis terhadap seluruh keberjalanan pengerjaan tugas akhir:
\begin{enumerate}[label=H\arabic*.]
    \item Skor pemahaman semantik hilir model \textit{brain-tuning} fMRI lebih tinggi dari model \textit{fine-tuning} bahasa saja dan model yang sinyal otaknya diacak (P1, P3).
    \begin{itemize}
        \item[$H_{0}$:] Skor pemahaman semantik hilir model dengan \textit{brain-tuning} tidak lebih tinggi daripada skor pemahaman semantik kedua model kontrol tersebut.
        \item[$H_{1}$:] Skor pemahaman semantik hilir model dengan \textit{brain-tuning} jelas lebih tinggi daripada skor pemahaman semantik kedua model kontrol tersebut.
    \end{itemize}
    \item Terdapat rentang bobot \textit{loss} sinyal otak yang teramati dalam menaikkan keselarasan otak tanpa penurunan berarti pada fungsionalitas model, sedangkan pada bobot \textit{loss} yang besar (di luar rentang) fungsionalitasnya menurun (P4).
    \begin{itemize}
        \item[$H_{0}$:] Tidak ditemukan rentang bobot \textit{loss} sinyal otak yang dapat menaikkan keselarasan otak tanpa penurunan fungsionalitas yang berarti. Fungsionalitas menurun di semua bobot, atau tidak menurun pada bobot besar.
        \item[$H_{1}$:] Ditemukan rentang bobot \textit{loss} sinyal otak yang dapat menaikkan keselarasan otak tanpa penurunan berarti pada fungsionalitas model, sedangkan pada bobot yang besar (di luar rentang) fungsionalitasnya menurun.
    \end{itemize}
    \item Perubahan representasi ($1-$CKA terhadap model awal) pada lapisan awal hingga tengah ($\le \ell^*$) lebih besar daripada pada lapisan di atas $\ell^*$, dan juga lebih besar daripada model yang diberi \textit{fine-tuning} pada bahasa saja (P2, P6).
    \begin{itemize}
        \item[$H_{0}$:] Perubahan pada lapisan awal hingga tengah tidak lebih besar daripada lapisan di atasnya atau juga tidak lebih besar daripada model bahasa saja.
        \item[$H_{1}$:] Perubahan pada lapisan awal hingga tengah jelas lebih besar daripada lapisan di atasnya dan juga jelas lebih besar daripada model bahasa saja.
    \end{itemize}
    \item Kenaikan keselarasan pada ROI (\textit{Region of Interest}) kategorikal seperti area otak tingkat tinggi terbukti lebih besar daripada area auditori primer otak (P2).
    \begin{itemize}
        \item[$H_{0}$:] Kenaikan keselarasan pada ROI (\textit{Region of Interest}) seperti area otak tingkat tinggi terbukti sama/lebih kecil daripada di area auditori primer.
        \item[$H_{1}$:] Kenaikan keselarasan pada ROI (\textit{Region of Interest}) seperti area otak tingkat tinggi jelas terbukti lebih besar daripada di area auditori primer.
    \end{itemize}
    \item Porsi keselarasan yang dijelaskan fitur teks tingkat rendah menurun setelah \textit{brain-tuning} pada fMRI asli dan tidak menurun pada fMRI yang diacak (P5).
    \begin{itemize}
        \item[$H_{0}$:] Porsi keselarasan dari fitur teks tingkat rendah tidak menurun setelah \textit{brain-tuning} fMRI asli, atau penurunannya sepadan dengan fMRI acak.
        \item[$H_{1}$:] Porsi keselarasan dari fitur teks tingkat rendah ikut menurun setelah \textit{brain-tuning} fMRI asli, serta penurunannya lebih besar dari fMRI acak.
    \end{itemize}
\end{enumerate}

Tiap hipotesis juga merupakan pemetaan dari setiap poin rumusan masalah di Subbab \ref{chap:pd_rm} yakni: (1) RM2 dipetakan pada H1 dan H2, (2) RM3 dipetakan pada H3, H4, dan H5, (3) RM1 tidak dipetakan pada H manapun karena bersifat perancangan.
\vfill


% \begin{enumerate}
%     \item \textbf{H1:} Model LLM hasil \textit{brain-tuning} memiliki kinerja pemahaman semantik yang lebih tinggi secara signifikan secara statistik dibandingkan model kontrol yang di-\textit{fine-tune} pada data yang sama tanpa sinyal otak dengan komputasi setara.
%     \begin{itemize}
%         \item $H_0$: tidak terdapat perbedaan kinerja pemahaman semantik antara model \textit{brain-tuning} dan model kontrol.
%         \item $H_1$: model \textit{brain-tuning} memiliki kinerja pemahaman semantik yang lebih tinggi daripada model kontrol.
%     \end{itemize}
%     \item \textbf{$H2:$} \textit{Brain-tuning} mengubah representasi model secara tidak merata antar lapisan, dengan perubahan terbesar terjadi pada lapisan tertentu yang selaras dengan hierarki pemrosesan bahasa di korteks.
%     \item \textbf{H3:} \textit{Brain-tuning} tidak menurunkan kapabilitas umum model secara signifikan pada tugas hilir di luar tugas pemahaman semantik utama.
% \end{enumerate}

% --- Metodologi ---
\section{Metodologi}
Tugas akhir ini menggunakan pendekatan penelitian eksperimental dengan desain kontrol yang ketat. Tahapan yang dilalui adalah sebagai berikut:
\begin{enumerate}
    \item \textbf{Investigasi latar belakang:} Mengumpulkan informasi dari beberapa riset relevan terkait \textit{brain-tuning} fMRI pada LLM yang berfungsi sebagai investigasi fakta kebaruan. Fakta tersebut direpresentasikan sebagai celah riset yang potensial untuk pengembangan lanjutan. Selain itu, fakta tersebut juga menjadi basis utama penyusunan rumusan masalah, tujuan, serta hipotesis tugas akhir. 
    \item \textbf{Studi literatur:} Menelaah sejumlah literatur mengenai sumber daya yang digunakan dalam pengerjaan tugas akhir. Sumber daya tersebut mencakup LLM sebagai model basis, fMRI naturalistik sebagai dataset utama, teknik \textit{fine-tuning} parameter efisien (PEFT) dengan LoRA, mekanisme pengaplikasian \textit{brain-tuning}, skor keselaran beserta representasi per lapisan model dan per ROI, serta \textit{benchmark} yang digunakan dalam mengevaluasi pemahaman semantik dengan tugas hilir. Selain itu, juga tersaji pengetahuan umum yang mendasari riset serta literatur terdahulu sebagai tolok ukur pengerjaan. Literatur diperoleh dari pencarian jurnal Q1 dan Q2 di Scopus dan sejumlah artikel dari berbagai sumber seperti ArXiv, IEEE Xplore, serta arsip terbuka lainnya.
    \item \textbf{Analisis kebutuhan dan pemilihan solusi:} Menentukan kebutuhan fungsional dan nonfungsional, alternatif dan pemilihan model, strategi \textit{fine-tuning}, dan objektif. Tahapan ini akan dipaparkan secara komprehensif pada Bab \ref{chap:analisis-masalah}.
    \item \textbf{Perancangan:} Membagi data latih dengan cerita yang tersedia, memilih data validasi penentu \textit{hyperparameter}, menguji data yang belum pernah digunakan, merancang objektif/\textit{loss} gabungan (\textit{loss} model dan otak), dan membangun tiga varian model: (1) \textit{brain-tuning} sebagai \textit{baseline}, serta \textit{fine-tuning} (2) bahasa saja dan (3) fMRI acak sebagai kontrol, dengan penyetaraan komputasi.
    \item \textbf{Implementasi:} Membangun \textit{pipeline} \textit{fine-tuning} LoRA (dan QLoRA secara opsional untuk model 4--8 miliar parameter) dengan pustaka PyTorch dan Hugging Face (\textit{Transformers} dan PEFT) pada platform Kaggle dengan GPU 16 GB (P100 atau T4). Bobot model diperoleh dari Kaggle Models, sedangkan konfigurasi, versi, dan \textit{random seed} dicatat untuk menjaga reprodusibilitas.
    \item \textbf{Evaluasi:} Membandingkan \textit{benchmark} pemahaman semantik (STS-B, WiC atau SuperGLUE, RTE atau MNLI, HellaSwag, MMLU, Perplexity), pengujian signifikansi dengan beberapa \textit{seed}, analisis per lapisan model dan per area otak, validasi lintas dataset (opsional), serta pengujian hipotesis (Subbab \ref{chap:pd_hipo}).
    \item \textbf{Penarikan kesimpulan dan saran:} Menyimpulkan seluruh hasil pengujian terhadap hipotesis serta merumuskan saran strategis untuk penelitian lanjutan.
\end{enumerate}
\vfill % memaksa konten di atas berada di bagian atas halaman

% --- Sistematika Penulisan ---
\section{Sistematika Penulisan}
Laporan tugas akhir ini disusun dengan sistematika penulisan sebagai berikut:
\begin{enumerate}
    \item \textbf{Bab I Pendahuluan:} Berisi latar belakang, rumusan masalah, tujuan, ruang lingkup, asumsi dan batasan, hipotesis, metodologi, serta sistematika laporan.
    \item \textbf{Bab II Studi Literatur:} Berisi kajian  mengenai pengetahuan kasus, landasan teori (LLM dan \textit{Transformer}, \textit{fine-tuning} dan LoRA, fMRI dan BOLD, \textit{encoding} model, \textit{brain-alignment} dan objektif, analisis per lapis, evaluasi semantik), serta penelitian terkait sebagai penegasan posisi riset tugas akhir ini.
    \item \textbf{Bab III Analisis Masalah:} Berisi analisis terhadap kondisi aktual, kebutuhan fungsional dan nonfungsional, alternatif solusi, dan pengukuhan solusi akhir.
    \item \textbf{Bab IV Desain Konsep Solusi:} Berisi konsep solusi terpilih, yaitu alur \textit{brain-tuning} (data, model, dan objektif gabungan \textit{loss}), desain kondisi kontrol, serta rancangan evaluasi dan analisis untuk membuktikan kebenaran tiap hipotesis.
    \item \textbf{Bab V Rencana Selanjutnya:} Berisi rencana pelaksanaan tugas akhir, tahapan implementasi dan evaluasi, jadwal pengerjaan, serta risiko dan mitigasinya.
    % \item \textbf{Bab V Implementasi}: Berisi proses implementasi kerangka \textit{fine-tuning} LoRA, teknologi yang digunakan, serta konfigurasi eksperimen.
    % \item \textbf{Bab IV Perancangan}: Berisi perancangan objektif gabungan \textit{loss} dari bahasa dan \textit{loss} dari penyelarasan otak manusia (\textit{brain-alignment}), serta desain kontrol eksperimen model \textit{brain-tuning}.
    % \item \textbf{Bab VI Evaluasi:} Berisi metode pengujian, hasil perbandingan model \textit{brain-tuning} dan kontrol, analisis representasi per-lapis, serta pembahasan terhadap hipotesis.
    % \item \textbf{Bab VII Kesimpulan dan Saran:} Berisi kesimpulan dari hasil tugas akhir dan saran untuk pengembangan lebih lanjut.
    \item \textbf{Daftar Pustaka:} Berisi daftar referensi yang dipakai guna menyusun laporan.
    \item \textbf{Lampiran:} Berisi dokumen pendukung yang tidak disertakan dalam naskah utama riset guna menjaga alur baca tanpa mengurangi kelengkapan informasi.
\end{enumerate}